% Generated draft: requires mathematical and source review.
\documentclass[11pt]{article}
\usepackage[margin=25mm]{geometry}
\usepackage{fontspec}
\setmainfont{Latin Modern Roman}
\usepackage{amsmath,amssymb,mathtools,unicode-math}
\setmathfont{Latin Modern Math}
\usepackage{enumitem,needspace,xurl,hyperref}
\setlength{\parindent}{0pt}
\setlength{\parskip}{5pt}
\setlength{\emergencystretch}{3em}
\newcommand{\N}{\mathbb{N}}
\newcommand{\Z}{\mathbb{Z}}
\newcommand{\Q}{\mathbb{Q}}
\newcommand{\R}{\mathbb{R}}
\newcommand{\C}{\mathbb{C}}
\newcommand{\F}{\mathbb{F}}
\newcommand{\A}{\mathbb{A}}
\newcommand{\PP}{\mathbb{P}}
\newcommand{\E}{\mathbb{E}}
\newcommand{\eps}{\varepsilon}
\newcommand{\dd}{\,\mathrm{d}}
\begin{document}
\section*{UnsolvedMath ID: 4000020}
\textbf{Source problem code:} catalogue entry
\subsection{Definitions and mathematical statement}
Use the standard simple-random-walk convention on a finite connected simple undirected graph $G=(V,E)$. Write $d_G$ for graph distance, $\deg(x)$ for degree, and
\[
m_x=\frac1{\deg(x)}\sum_{y\sim x}\delta_y,\qquad
(Pf)(x)=\frac1{\deg(x)}\sum_{y\sim x}f(y).
\]
With $W_1$ computed from $d_G$, put $\kappa_G(x,y)=1-W_1(m_x,m_y)/d_G(x,y)$ for $x\ne y$. The spectral gap is $1-\lambda_2(P)$, where eigenvalues are ordered by their real value, equivalently the first positive eigenvalue of $I-P$ for the reversible walk.

Does there exist a sequence $(G_j)_{j\geq1}$ and constants $D<\infty$, $\varepsilon>0$ such that, for every $j$,
\[
\max_{x\in V(G_j)}\deg(x)\leq D,\qquad
1-\lambda_2(P_j)\geq\varepsilon,\qquad
\kappa_{G_j}(x,y)\geq0\quad(x\ne y),
\]
and $\operatorname{diam}(G_j)\to\infty$? This states bounded-degree expansion and nonnegative coarse Ricci curvature with a common graph-walk convention.

The supplied literature assessment reports a negative answer by Salez (2022). It additionally reports a stronger obstruction allowing vanishing proportions of large degrees, large eigenvalues, and negatively curved edges, and an analogous obstruction under Bakry--\'Emery $CD(0,\infty)$ curvature. Those refinements are source reports; their full hypotheses and quantifiers are not included and are not asserted here as a separate formal theorem. The original item does not explicitly fix walk laziness or spectral normalization, so these conventions must be matched when applying the cited result.
\subsection{Short English statement}
Can bounded-degree graphs form an expander family of growing diameter while having nonnegative transport-based Ricci curvature? The supplied Salez reference reports that no such family exists.
\subsection{Sources}
\begin{itemize}
\item Yann Ollivier, \emph{Discrete Ricci curvature: Open problems} (2008), Problem T, \url{https://www.yann-ollivier.org/rech/publs/problems_curvmarkov.pdf}; accessed 2026-07-29, with PDF-text extraction recorded in the supplied provenance.
\item Justin Salez, ``Sparse expanders have negative curvature,'' \emph{Geometric and Functional Analysis} 32 (2022), 917--948, DOI \url{https://doi.org/10.1007/s00039-022-00618-3}, arXiv:2101.08242. The supplied summary attributes the original problem to A. Naor and E. Milman, publicized by Y. Ollivier (2010), and the related $CD(0,\infty)$ conjecture to Cushing--Liu--Peyerimhoff.
\item Local registry \texttt{data/unsolvedmath/problems.json}, permanent UnsolvedMath ID $4000020$, source code catalogue entry, \emph{Expanders}. Dataset attribution: Ulam AI and the UnsolvedMath Contributors.
\end{itemize}
\end{document}
